\properlane{eschatological-detachment}{The World Grown Old and the Coming of the Strict Judge}
\runninghead{The World Grown Old and the Strict Judge}

The Gospel of the First Sunday of Advent confronts the listener with two abrupt reversals inside six sentences. First, Christ foretells signs in the sun, moon, and stars, nations trapped in distress at the roaring sea, men withering away in terror, and the powers of heaven shaken, and then turns immediately to His disciples with a command of joy: ``when these things begin to come to pass, look up, and lift up your heads, because your redemption is at hand'' (Luke~21:28). Second, He illustrates this cosmic upheaval with trees putting forth their fruit before summer (Luke~21:29--30) and contrasts the passing of sky and earth with the permanence of His own spoken words (Luke~21:33). Why should collapse be a reason to rejoice, and what kind of fruit can a dying world possibly produce?

Seven Fathers and Doctors expound these passages and related Scriptures in the light of the Last Judgment. St.~Gregory the Great expounds Luke~21:25--33 as Christ's warning to detach our hearts from an aging world before the strict Judge; St.~Cyril of Alexandria explains Luke~21:25--28 as the promise of bodily resurrection and cosmic refashioning; St.~Bede expounds Luke~21:28--33 on how heaven and earth pass away in outward form while remaining in substance, applying to verse~33 the arguments St.~Augustine develops in \work{De civitate Dei}~XX on cosmic transmutation, renewal through fire, and the lower atmospheric heavens; St.~John Chrysostom preaches Romans~13:11--14 as a pre-dawn summons to vigilance before the Judgment; St.~Leo the Great applies Luke~21:34 directly to Advent sobriety; and St.~Robert Bellarmine expounds Psalm~24:1--3 as tearing the soul away from earthly briers. Read together across the 1962 Mass, their expositions show the liturgical year opening as a night watch before the strict Judge, whose appearing terrifies lovers of this passing age and brings final deliverance to those who wait on Him in detachment.

\subsection{Why Cosmic Ruin Commands Joy}

St.~Gregory the Great (\work{Homiliae in Evangelia}~1.1--6) takes up the Gospel's reversal directly: ``Our Lord and Saviour wisheth to find us ready at His second coming. Therefore He telleth us what will be the evils of the world as it groweth old, that He may wean our hearts from worldly affections'' (\S1; Bute~1908). When Gregory asks what Christ means by \latin{virtutes caelorum} (``the powers of heaven'') being moved in verse~26, he identifies them as the angelic orders---angels, archangels, thrones, dominations, principalities, and powers---who will appear visibly at the coming of the strict Judge to demand an exact reckoning of what our invisible Creator now bears with patience (\S2). Those who refused to hear Him in humility will then behold Him in power and majesty. In the Mass of this Sunday, the \latin{Credo} confesses that same Judge returning in glory to judge the living and the dead, while the Preface of the Most Holy Trinity joins the congregation to the unceasing praise of angels, archangels, cherubim, and seraphim.

Only after facing that reckoning does Gregory turn to verse~28 (\S3; Oxford~1843, preserved in St.~Thomas Aquinas, \work{Catena Aurea in Lucam} on Luke~21:28). In Scripture, ``head'' often stands for the mind, because as the body's members are ruled by the head, human thoughts are directed by the mind. To lift up your heads (\latin{levate capita}) therefore means to gladden your hearts (\latin{exhilarate corda}) and raise the mind toward the joys of the heavenly country. Gregory draws the dividing line from James~4:4: whoever wills to be a friend of this world becomes an enemy of God. Anyone who grieves to see this passing age approach its end betrays where his heart is rooted. Mortal life is a road, not a home; Gregory asks his Roman listeners how strange it is to grow exhausted by the toil of a journey and still wish that the road would never finish.

From that dividing line Gregory draws his severe reading of the fig tree and all the trees in verses~29--30 (\S3). Where a hearer expects ``fruit'' to mean flourishing health, Gregory turns the parable inside out: \latin{fructus mundi ruina est}---``the fruit of the world is its ruin.'' The world grows only so that it may fall; it sprouts only so that disasters may consume whatever it puts forth. Just as ripening fruit tells a farmer that summer is close, the world's visible decay tells a believer that the eternal summer of God's kingdom is near, when the clouds of earthly sorrow will pass away and the days of life will shine with the brightness of the eternal Sun.

Gregory presses the warning from a local catastrophe that had struck Rome two days before he preached (\S\S5--6). Comparing the aging world (\latin{senescentem mundum}) to an old man whose back bows, whose neck dries, and whose panting cuts off his speech, he recalls the sudden whirlwind (\latin{nudius tertius}, \latin{subito turbine}) that tore ancient trees out by the roots, wrecked houses, and threw down churches from their foundations, catching men in sudden death on the night they were planning the next morning's business (Bl.~Ildefonso Schuster, I, p.~322, inadvertently remembers this windstorm as an earthquake). If the invisible Judge can overturn stone buildings by stirring a single breath of wind and a thin cloud, Gregory asks, what will happen when He appears in person? Present calamities are only heralds running ahead of coming wrath (\latin{sequentis irae praecones}, \S6), warning believers to amend their lives before the sentence falls.

\subsection{Resurrection, Renewal, and the Abiding Word}

Two patristic arguments keep Gregory's renunciation of the world from sliding into contempt for God's creation. First, St.~Cyril of Alexandria (\work{Commentary on Luke}, Sermon~CXXXIX, Payne Smith pp.~653--654; leaving vv.~29--33 unexpounded) marks verse~25 as the Lord's shift from the siege of Jerusalem to the end of the age, and explains what verse~28's ``your redemption is near'' promises: the resurrection of the dead and the refashioning of creation to what it was at the beginning. When Christ comes openly as God and Lord, ``this earthly and infirm body shall put off corruption, and shall clothe itself with incorruption by Christ's gift,'' conformed to His glorious body. Joining 2~Peter~3:10--11 with Luke~21:36, Cyril draws the direct moral consequence of watchful holiness before the judgment seat.

Second, St.~Bede (\work{In Lucae Evangelium Expositio}~VI, on vv.~28--29, 32--33) takes up verses~32--33 directly. After transmitting Gregory's exposition of verses~28--29, Bede pauses on \latin{Amen} (noting that neither Greek nor Latin translators dared render the Hebrew word for ``truly'' into common speech lest its solemnity grow cheap, \latin{ne vilesceret nudatum}) and explains \latin{generatio haec} (``this generation'') as either the whole human race or specifically the Jewish people. On verse~33---``Heaven and earth shall pass away, but my words shall not pass away''---Gregory (\S4) had already pointed out the double paradox: nothing in bodily nature appears sturdier than sky and earth, and nothing passes faster than a spoken utterance, which finishes only by ceasing to sound; yet Christ's spoken word expresses unchanging truth while the physical frame passes. Bede then answers two difficulties raised by verse~33. First, citing 2~Peter~3:5--7, he distinguishes the upper starry heaven (\latin{aethereum, sive sidereum}) from the lower airy heaven (\latin{aereum}) where birds fly and clouds gather: because the waters of the Flood rose fifteen cubits above the mountaintops without reaching the starry firmament, judgment's fire will sweep only that same lower atmospheric region. Second, reconciling Luke~21:33 with Ecclesiastes~1:4 (``the earth standeth for ever'') through 1~Corinthians~7:31, Apocalypse~21:1, and Psalm~101:27, Bede explains that heaven and earth pass away in their present outward form (\latin{per imaginem}), yet endure without end in their substance (\latin{per essentiam}). Cleansed of its present appearance by fire (\latin{per ignem tergetur}), creation is preserved in its own nature (\latin{in sua semper natura servabitur}), just as bare winter ground turns green in spring and the night sky is renewed each morning by daylight.

Bede applies to Luke~21:33 three distinct arguments that St.~Augustine develops in Book~XX of \work{De civitate Dei}, though Augustine himself does not expound Luke~21:33 directly. In \work{De civitate Dei}~XX.14---expounding Apocalypse~20:9--13 (especially verse~11, where heaven and earth flee from the throne), looking ahead to Apocalypse~21:1, and quoting 1~Corinthians~7:31--32---Augustine teaches that ``this world shall pass away by transmutation, not by absolute destruction'': ``The figure, therefore, passes away, not the nature.'' In \work{De civitate Dei}~XX.16, expounding Apocalypse~21:1 (``a new heaven and a new earth'') alongside Apocalypse~20:11 and 1~Corinthians~7:31, he explains how that renewal comes about: ``then shall the figure of this world pass away in a conflagration of universal fire,'' just as the world was once flooded by water, so that ``the world itself is renewed to some better thing,'' fittingly accommodated to men ``themselves renewed in their flesh to some better thing.'' Finally, in \work{De civitate Dei}~XX.24, Augustine answers the pagan philosopher Porphyry, who praised Hebrew piety while accusing Christians of folly for expecting the world to perish. Appealing to Psalm~101:26--28 (``The heavens are the work of Thy hands; they shall perish''), Augustine notes in passing that 1~Corinthians~7:31, 1~John~2:17, and Matthew~24:35 (not Luke~21:33) use milder expressions than ``They shall perish,'' and correlates 2~Peter~3:6--11 with the Flood to argue that only the stormy lower atmospheric heavens perish while the higher heavens in whose firmament the stars are set remain safe.

\subsection{Waking at Dawn and Waiting Without Shame}

St.~John Chrysostom (\work{Homilies on Romans}~24) preaches St.~Paul's summons in Romans~13:11--14 in view of the Resurrection and Judgment. Why does Paul tell the Roman Christians that their salvation is nearer now than when they first believed? Chrysostom notices a familiar human drift: beginners start with full ardor and slacken as years pass. Paul demands the opposite, comparing believers to racers in a stadium who push harder when they sight the end of the course and the prize. As time runs short and the reckoning draws closer, listlessness must yield to keen readiness.

Chrysostom illustrates ``the night is far spent, the day is at hand'' (Rom.~13:12) from ordinary morning life: when people hear the swallow twittering before dawn, neighbors wake each other while it is still dark, shaking off dreams and getting dressed so sunrise will not catch them still yawning in bed. Then Chrysostom pauses over a verbal shift in Paul's Greek and Latin: Paul calls sins ``works of darkness'' (\latin{opera tenebrarum}), yet calls virtues not ``works'' but ``armour of light'' (\latin{arma lucis}). Ordinary steel armor is heavy and grim to put on; Paul's armor is light itself, guarding the Christian while making him shine brighter than a sunbeam. Under such a Commander, spiritual warfare becomes a solemn festival, and the believer steps out ``at once soldier and bridegroom.'' Paul then speaks as though it were already daylight (v.~13), gently saying ``let us walk'' while condemning drunken revelry, chambering, and faction---which Cornelius a Lapide (\work{Comm.~in~omnes~D.~Pauli~Epistolas}, pp.~196--197) explains from the Greek \emph{komos} (nocturnal carousing, p.~197) after treating \latin{arma lucis} as offensive and defensive warfare against the world, the flesh, and the devil (p.~196)---before ascending in verse~14 from works to armor to putting on the King Himself, girt about the believer while dwelling within as Way, Bridegroom, Food, High Priest, and Friend. St.~Leo the Great (\work{Sermon}~19.1), preaching on Luke~21:34 immediately after this Gospel, draws the corresponding Advent discipline: fasting, prayer, and almsgiving keep the heart from being weighed down by gluttony and worldly cares before that sudden Day.

The psalm chants and opening prayers voice this same watchful detachment. At the Introit (\latin{Ad te levavi}, Ps.~24:1--3, 4b), St.~Robert Bellarmine (\work{Explanatio} on Ps.~24:1--3) hears a soul that has found only ``briers and thorns'' everywhere in creatures, tearing itself free from earthly attachments to entrust itself to God alone against visible and demonic foes, encouraged by the patriarchs' assurance that none who wait in sure hope are put to shame. Repeating that antiphon verbatim at the Offertory and asking in the Gradual for God to make His ways known (\latin{notas fac mihi}, Ps.~24:4b) sets the congregation's trust before the altar, while the Collect (\latin{Excita, quaesumus}) names the peril already pressing upon sinners---\latin{ab imminentibus peccatorum nostrorum periculis}, ``approaching dangers brought on by our sins''---and asks Christ to rouse His power to snatch (\latin{eripi}) and save (\latin{salvari}) us. At the altar, the Secret (\latin{Haec sacra nos}) asks the Mysteries to cleanse worshippers so they may stand purer before their divine Source (\latin{principium}).

William Durandus (\work{Rationale}~VI.2--3, fols.~256v, 261v) pairs violet restraint and the silence of \latin{Gloria in excelsis} with the Sunday Alleluia (Ps.~84:8, pleading for God's mercy and saving gift) as the voice of a season ``partly of rejoicing and partly of mourning'' (\latin{partim exultationis, et partim moeroris}), and reads \latin{Ad te levavi}, \latin{propior est nostra salus}, and \latin{Erunt signa} together as announcing the nearness of God's kingdom. In the same sober watchfulness, the congregation sings the Communion antiphon (Ps.~84:13) trusting God's goodness (\latin{benignitas}) to bring forth faithful fruit, and prays in the Postcommunion (\latin{Suscipiamus}) with the approaching feast in view alongside judgment.

\subsection{Tensions and Limits of This Perspective}

Even within this group of Fathers, a real tension separates Gregory's stark aphorism ``the fruit of the world is its ruin'' from the theology of cosmic restoration taught by Cyril of Alexandria, Bede, and Augustine. Taken by itself, Gregory's portrait of a decrepit, collapsing world sounds darker than the whole scriptural picture unless read beside Cyril's promise of bodily incorruption and Bede's teaching that heaven and earth are purged and renewed rather than annihilated. Moreover, interpreting the Mass primarily in the light of the Last Judgment accounts less directly for Rome's choice of the Marian basilica of St.~Mary Major (\latin{ad praesepe}, ``at the Crib'') for this Sunday's station, and leaves to one side the Incarnational and Marian expositions of Luke~21:27 and Psalm~84:13 developed by St.~Ambrose and St.~Robert Bellarmine.

\subsection{The Four Senses}

\begin{fourSenses}
\item[Literal.] Psalm~24 voices a prayer of unashamed trust in God and asks instruction in His ways against mocking foes; Psalm~84 begs God to restore mercy and fruitfulness to His people; St.~Paul urges Roman Christians to wake from sleep, abandon nocturnal vice, and put on Christ because salvation draws nearer; in the Temple courts during Passion Week, Jesus foretells cosmic signs, the coming of the Son of Man, the budding of the fig tree and all trees, and the permanence of His words (see \hyperref[sec:date-location]{\work{Scriptural Date and Location}} for scriptural chronology).
\item[Allegorical.] The shaken \latin{virtutes caelorum} signify the angelic choirs appearing as ministers of the strict Judge (Gregory, Hom.~1.2); the \latin{arma lucis} signify Christian virtues shining as festive battle-dress in Christ the King (Chrysostom, Hom.~24); the fig tree's ``fruit'' signifies the world's own ruin signaling the approach of God's kingdom (Gregory, Hom.~1.3).
\item[Moral.] Detachment from a passing age and its briers (Gregory, Bellarmine); sober vigilance through prayer, fasting, and almsgiving against gluttony, impurity, and faction (Rom.~13:13; Leo, Serm.~19.1); patient, unashamed perseverance (\latin{exspectatio}) in God's ways during the night watch.
\item[Anagogical.] The public Parousia of the Son of Man in power and majesty; the cleansing and renewal of heaven and earth in outward form (\latin{per imaginem}) rather than destruction in substance (\latin{per essentiam}) through judgment's fire (Bede on Luke~21:33, \latin{per ignem tergetur}; Augustine, \work{De~civ.~Dei}~XX.14, 16 on transmutation and renewal through universal fire); the bodily resurrection of the saints into incorruption (Cyril, Serm.~CXXXIX); and the eternal summer of God's kingdom (Gregory, Hom.~1.3).
\end{fourSenses}
